\section{Verification purpose}

\textbf{Case ID:} \texttt{\CaseID}. This case verifies raster lookup and fire-behavior coupling when a propagating front crosses four fuel-model regions.  Flat terrain and zero wind isolate fuel identity.  The verification question is whether every quadrant's maximum spread rate is within 10 percent of its fuel-specific reference.

\section{Mathematical formulation and numerical implementation}

The level-set normal speed is evaluated from the fuel model at each active cell.  Because the environmental forcing is uniform, changes in \(V_n\) across the domain should follow only changes in the \texttt{fbfm} raster.  Correct row orientation and quadrant masking are therefore part of the test.

\section{Derivation of expected behavior}

The independent regional targets are \(0.080\), \(0.472\), \(1.492\), and \(0.871~\mathrm{m\,min^{-1}}\) for Fuel Models 8, 7, 4, and 2 in the northwest, northeast, southwest, and southeast quadrants, respectively.  These values also predict a visibly noncircular front after it enters all four regions.

\section{Verification metrics and acceptance criteria}

The maximum finite \texttt{vs} in each quadrant is compared with its corresponding target.  Regional maxima preserve the local head-fire response and avoid mixing physically distinct fuels.  Every component uses 10 percent relative tolerance, and PASS requires all four components.

\section{Simulation configuration and predicted outcome}

Four equal quadrants on a \(126\times126\), 30 m grid contain Fuel Models 8, 7, 4, and 2. Terrain is flat, wind is zero, moisture is uniform, and ignition is central. The fastest extension is predicted in Fuel Model 4 and the slowest in Fuel Model 8.

\section{Actual simulation results}

The required evidence is generated only from case-local ELMFIRE outputs. The current postprocessing status is \textbf{\VerificationStatus}; \PassingMetricCount{} of \MetricCount{} required metric rows pass across \VariantCount{} variant(s). A missing output is reported as NOT RUN or NOT EVALUABLE and is not interpreted as verification evidence.


\begin{figure}[H]
  \centering
  \IfFileExists{../figures/whole_domain_result.pdf}{%
    \includegraphics[width=0.82\textwidth,height=0.58\textheight,keepaspectratio]{whole_domain_result.pdf}%
  }{\fbox{\parbox{0.85\textwidth}{No case-local whole-domain ELMFIRE raster is available.}}}
  \caption{Representative whole-domain case-local ELMFIRE output. The generated figure records its source path, coordinate reference system, and variant.}
\end{figure}


\begin{figure}[H]
  \centering
  \IfFileExists{../figures/verification_summary.pdf}{%
    \includegraphics[width=\textwidth,height=0.58\textheight,keepaspectratio]{verification_summary.pdf}%
  }{\fbox{\parbox{0.85\textwidth}{No case-local summary figure is available.}}}
  \caption{Case-local expected and measured verification metrics. A missing measured value denotes unavailable or incomplete ELMFIRE output.}
\end{figure}

\section{Calculated metrics and verification decision}

\begin{center}
\begin{tabular}{@{}llll@{}}
\toprule
Variant / metric & Expected & Calculated & Status \\
\midrule
\MetricRows
\bottomrule
\end{tabular}
\end{center}

The authoritative overall decision is \textbf{\VerificationStatus} and is generated by \texttt{scripts/postprocess.py} from \texttt{outputs/metrics.json}; it is not recomputed in this report.

\section{Technical interpretation}

A quadrant swap indicates raster orientation or mask errors; a fuel-specific discrepancy indicates fuel-table selection or fire-behavior coupling. No scientific conclusion should be drawn when the status is NOT RUN or NOT EVALUABLE.

\section{Scope and limitations}

The test samples four fuel models under one moisture state and does not validate the empirical fuel parameters. This case provides implementation-verification evidence only.
